\chapter{Procedimiento metodológico}
\label{chapter:chapter04}

Este capítulo explica cómo se pasó del problema observado en campo a una metodología de captura, análisis e integración de la solución perimetral. No se parte de un escenario ideal: la finca opera con restricciones eléctricas, administrativas y de red que condicionan la trazabilidad de la telemetría durante el ordeño. Por eso, el procedimiento se organiza en cinco momentos: reconocimiento del entorno, obtención y análisis de la información, evaluación de alternativas, integración de la solución y reevaluación previa a la validación final.


\section{Reconocimiento metodológico del entorno}
\label{sec:reconocimiento-definicion-problema}
El problema de investigación fue planteado en el Capítulo~\ref{chapter:chapter01}. Aquí se retoma desde el trabajo de campo: se explica cómo se reconoció en la finca, qué evidencia permitió delimitarlo y cómo esas observaciones se convirtieron en variables, muestras y criterios de integración de la solución.

\subsection{Criterios de observación del entorno}
La etapa de reconocimiento se organizó en tres frentes: inspección física y eléctrica del entorno de ordeño, levantamiento de la topología operativa del ecosistema comercial \textit{SenseHub Dairy} y captura pasiva de los dominios serial y de red. Esta división ayudó a separar lo que pertenece al contexto productivo, lo que depende de las restricciones del sistema comercial y lo que sí podía medirse de forma reproducible.

En el frente físico se registraron condiciones capaces de afectar la continuidad del equipo: humedad relativa superior al 80\,\%, variaciones de temperatura, alta incidencia de descargas atmosféricas, inestabilidad del suministro eléctrico y deficiencias de equipotencialidad. Estas observaciones justifican una intervención no invasiva y orientan las mediciones de diagnóstico eléctrico.

En el frente de red y comunicación se documentaron interrupciones intermitentes, variación temporal de la latencia, tráfico administrativo y comunicaciones hacia destinos externos que podían mezclarse con el flujo biótico. De acuerdo con los parámetros de la Unión Internacional de Telecomunicaciones (\acs{UIT}) \cite{itu_iot_latency} y los lineamientos del Grupo de Trabajo de Ingeniería de Internet (\acs{IETF}) \cite{rfc3393}, estas condiciones orientaron la selección de métricas asociadas con variación temporal de la latencia, disponibilidad del flujo, exposición de tráfico y capacidad de almacenamiento local con reenvío.

\subsection{Levantamiento topológico y puntos de captura}
La topología general presentada en la Figura~\ref{fig:contexto_general_cap1} se verificó mediante observación directa durante las visitas técnicas. En la finca, la recepción de señales se coordina mediante controladores \textbf{DF-800}, enlazados por un bus diferencial \textbf{RS-485} y acoplados a la computadora de control mediante \textbf{RS-232} a 19\,200\,bps (8N1), según la documentación del fabricante \cite{allflex2021,sensehub2023}.

El levantamiento permitió ubicar seis puntos de interés: collares, antena, controlador maestro, enlace serial, gabinete de red y computadora de la lechería. La inspección fotográfica confirmó dónde era posible observar sin intervenir, como se resume en la Figura~\ref{fig:foto_flujo}. A partir de ahí, el análisis se concentró en mediciones reproducibles del enlace serial, la \ac{LAN} y las capturas \ac{PCAP}.

\begin{figure}[H]
    \centering
    \begin{subfigure}[t]{0.32\linewidth}
        \centering
        \includegraphics[width=\linewidth,height=0.16\textheight,keepaspectratio]{Figuras/vacas.png}
        \caption{Collares de monitoreo en el hato.}
    \end{subfigure}\hfill
    \begin{subfigure}[t]{0.32\linewidth}
        \centering
        \includegraphics[width=\linewidth,height=0.16\textheight,keepaspectratio]{Figuras/antena.png}
        \caption{Antena de cobertura del sistema.}
    \end{subfigure}\hfill
    \begin{subfigure}[t]{0.32\linewidth}
        \centering
        \includegraphics[width=\linewidth,height=0.16\textheight,keepaspectratio]{Figuras/caja_master.png}
        \caption{Controlador maestro DF-800.}
    \end{subfigure}

    \vspace{0.6em}
    \begin{subfigure}[t]{0.32\linewidth}
        \centering
        \includegraphics[width=\linewidth,height=0.16\textheight,keepaspectratio]{Figuras/cable_serial.png}
        \caption{Enlace serial de captura.}
    \end{subfigure}\hfill
    \begin{subfigure}[t]{0.32\linewidth}
        \centering
        \includegraphics[width=\linewidth,height=0.16\textheight,keepaspectratio]{Figuras/caja_red.png}
        \caption{Gabinete de red local.}
    \end{subfigure}\hfill
    \begin{subfigure}[t]{0.32\linewidth}
        \centering
        \includegraphics[width=\linewidth,height=0.16\textheight,keepaspectratio]{Figuras/pc_lecheria.png}
        \caption{Computadora de la lechería.}
    \end{subfigure}
    \caption{Componentes del entorno de red, condiciones operativas y nodos de captura para el flujo de telemetría \textit{SenseHub Dairy} en la Finca La Esmeralda (elaboración propia).}
    \label{fig:foto_flujo}
\end{figure}

La observación de campo confirmó que la identidad visible durante el ordeño corresponde al identificador del collar Allflex, mientras que la identidad oficial de trazabilidad puede mantenerse en registros externos asociados al bolo intraruminal HDX u otros sistemas regulatorios \cite{datamars_hdx,trazar_agro,senasa_sira}. Esta distinción evitó tratar el identificador capturado como si fuera, por sí solo, la identidad oficial completa del animal.

\subsection{Caracterización temporal, eléctrica y de red}
A partir del problema delimitado en el Capítulo~\ref{chapter:chapter01}, la fase de reconocimiento caracterizó la sincronización entre activación de fotocelda, lectura del identificador del collar y llegada de telemetría biótica al nodo central. Durante los turnos de ordeño se registraron eventos físicos y tráfico asociado para evaluar si la evidencia conservaba una relación temporal suficiente con el evento productivo.

En el plano temporal y de red, las primeras capturas documentaron picos de variación temporal de la latencia de 71,3\,ms y tráfico de multidifusión a 7,14\,Hz. Estos valores sirvieron para fijar métricas posteriores de línea base, exposición de tráfico, eficiencia de extracción y concordancia entre evidencia serial y de red.

En el plano eléctrico, el análisis osciloscópico entre el chasis de la computadora y la tierra del bus de campo evidenció una señal parásita de 60,98\,Hz con una amplitud de 1,02\,Vpp, como se muestra en la Figura~\ref{fig:hantek_ground_loop}. La medición de referencia entre tierras registró, además, una resistencia de $3{,}37\,\text{k}\Omega$, valor que confirmó un acoplamiento eléctrico no despreciable entre dominios. Estas mediciones justificaron el criterio de aislamiento galvánico y alta impedancia de entrada para cualquier acople sobre el enlace serie.

\begin{figure}[H]
    \centering
    \includegraphics[width=0.8\linewidth,keepaspectratio]{Figuras/hantek_ground_loop.png}
    \caption{Evidencia de bucle de tierra: Componente de 60\,Hz detectada entre la tierra del bus y el chasis del nodo central (PC), evidenciando un diferencial de 1,02\,Vpp (elaboración propia).}
    \label{fig:hantek_ground_loop}
\end{figure}

\subsection{Restricciones operativas de intervención}
El reconocimiento de campo también dejó claras las restricciones prácticas de la intervención. El personal del \ac{PPA} no dispone de permisos administrativos sobre el nodo central; su acceso se limita a un perfil de usuario estándar, condición verificada mediante observación directa y entrevista con el coordinador técnico del programa. Esta limitación impide instalar herramientas de diagnóstico o crear almacenamiento temporal interno en la computadora de la lechería.

Las pruebas de campo también mostraron que la conexión física directa al bus \ac{RS-232} podía perturbar la estabilidad del enlace serie. Se detectaron diferencias de potencial de hasta 59,2\,V entre referencias de masa, por lo que la metodología descartó una intervención invasiva y adoptó como criterio de integración una captura perimetral, aislada y reversible.

\subsection{Matriz de evidencia diagnóstica}
La Tabla~\ref{tab:matriz_problematica} resume los diez factores críticos usados para traducir el problema del Capítulo~\ref{chapter:chapter01} en variables observables, criterios de diseño y criterios de validación. La matriz no repite la justificación inicial; ordena la evidencia por factor, valor observado y uso metodológico.

\begingroup
\scriptsize
\setlength{\tabcolsep}{2pt}
\renewcommand{\arraystretch}{1.08}
\begin{longtable}{|>{\raggedright\arraybackslash}p{2.85cm}|>{\raggedright\arraybackslash}p{3.60cm}|>{\raggedright\arraybackslash}p{6.60cm}|>{\raggedright\arraybackslash}p{2.75cm}|}
\caption{Matriz diagnóstica de factores críticos usados para operacionalizar el problema.}
\label{tab:matriz_problematica}\\
\hline
\rowcolor{gray!20} \centering\arraybackslash\textbf{Factor crítico} & \centering\arraybackslash\textbf{Evidencia o valor observado} & \centering\arraybackslash\textbf{Criterio metodológico derivado} & \centering\arraybackslash\textbf{Fuente} \\ \hline
\endfirsthead
\endhead
\hline
\endfoot
\textbf{Seguridad eléctrica} & Diferencial de potencial de hasta 59,2\,V; componente de 60,98\,Hz con 1,02\,Vpp; resistencia de $3{,}37\,\text{k}\Omega$ entre tierras. & Incorporar aislamiento galvánico, alta impedancia de entrada y descarte de intervención directa sobre el enlace serie. & Osciloscopio Hantek 2D72 y medición eléctrica de campo \\ \hline
\textbf{Ruido e integridad serial} & Transitorios de 885\,mVpp y perturbaciones cercanas a 4\,MHz en la línea de recepción. & Usar captura pasiva, validar errores de trama y conservar evidencia física complementaria. & Análisis de transitorios \\ \hline
\textbf{Variación temporal de latencia} & Picos de variación temporal de 71,3\,ms y línea base temprana de 61,250\,ms. & Definir ventanas de captura, métricas de línea base y tolerancias para correlación serial--red. & Trazas \acs{PCAP} \\ \hline
\textbf{Multidifusión y ruido de red} & Inundación de multidifusión a 7,14\,Hz; en el conjunto de diagnóstico, tráfico de multidifusión de hasta 18,99\,\%. & Segmentar el dominio de red, controlar ráfagas y separar tráfico administrativo del flujo biótico. & Diagnóstico de red\\ \hline
\textbf{Inestabilidad \acs{ARP}} & Conflictos \acs{ARP}: 484 casos en 62 visitas; una \acs{MAC} respondió por varias \acs{IP} en 9040 casos. & Aislar el segmento perimetral y usar conmutación controlada para reducir ambigüedad de red. & Capturas forenses \acs{PCAP} \\ \hline
\textbf{Exposición de información} & Tráfico en claro en la \ac{LAN}; telemetría \ac{UDP} asociada al puerto~6001 y tráfico HTTP observado. & Evaluar exposición pasiva y exigir transporte protegido, lista de destinos permitidos y publicación segura. & Diagnóstico de red\\ \hline
\textbf{Soberanía operativa} & Acceso restringido a la computadora de la lechería y ausencia de permisos administrativos. & Diseñar una pasarela perimetral independiente del nodo propietario, con almacenamiento local y reenvío. & Observación de campo \\ \hline
\textbf{Observabilidad local} & Ausencia de una vista local integrada para revisar alertas, sesiones, cobertura serial y cobertura \acs{PCAP}. & Incorporar paneles de consulta y salidas auditables que permitan seguimiento local sin depender del nodo propietario. & Observación de campo y panel Aletheia \\ \hline
\textbf{Continuidad de publicación} & Conectividad rural variable y dependencia de destinos externos para la publicación de evidencia normalizada. & Usar cola local, operación residente y mecanismo de almacenamiento y reenvío para evitar pérdida silenciosa de lotes. & Bitácora operativa y pruebas de pasarela \\ \hline
\textbf{Trazabilidad evento--animal} & Eventos de vaca sin identificador de collar asociado: 79 casos en 44 visitas; sesiones sin pareos serial--red. & Priorizar correlación evento--campo, conservar casos con $\eta=0\%$ y fijar criterios de inclusión muestral. & Serial, \acs{PCAP} y bitácora \\ \hline
\end{longtable}
\endgroup



La evidencia de ruido transitorio se conserva como registro complementario en la Figura~\ref{fig:hantek_noise}, porque permite contrastar visualmente la degradación física del enlace con las métricas de captura usadas en las etapas posteriores.

\begin{figure}[H]
    \centering
    \includegraphics[width=0.8\linewidth,keepaspectratio]{Figuras/hantek_noise.png}
    \caption{Interferencia electromagnética transitoria en la línea de datos, con picos de 885\,mVpp (elaboración propia).}
    \label{fig:hantek_noise}
\end{figure}

\section{Obtención, procesamiento y análisis de la información}
Con el problema reconocido, la investigación avanzó hacia la obtención y el tratamiento de evidencia verificable. Esta fase articuló el enfoque de investigación, la definición de variables, la estrategia de muestreo, el procesamiento de capturas y la validación del instrumento de análisis.

\subsection{Enfoque y diseño de la investigación}
Delimitado empíricamente el problema, la investigación adopta un \textbf{enfoque cuantitativo}, centrado en recolectar y analizar datos numéricos sobre el flujo telemétrico de la Finca La Esmeralda. Siguiendo a Montgomery \cite{montgomery2019}, este enfoque permite trabajar con magnitudes observables y comparables: latencia, variación temporal de la latencia, eficiencia de extracción ($\eta$), exposición de tráfico en claro y tiempo medio de recuperación del sistema.

El diseño de investigación es \textbf{cuasi-experimental aplicado}, con comparación antes/después de la integración perimetral. Esta formulación se ajusta al contexto real de la finca: no se aleatorizan animales ni se detiene la operación productiva, sino que se fija una condición previa y una condición posterior bajo criterios de captura comparables. Las variables de contraste son cobertura efectiva de captura, concordancia evento--campo, exposición de telemetría en claro y resiliencia operativa.

\subsection{Operacionalización de variables}
Para asegurar que los objetivos específicos se traduzcan en métricas verificables y comparables, se definieron las variables de estudio presentadas en la Tabla~\ref{tab:operacionalizacion}.

\begin{table}[!htbp]
\centering
\caption{Cuadro de operacionalización de variables de estudio.}
\label{tab:operacionalizacion}
\scriptsize
\setlength{\tabcolsep}{2pt}
\renewcommand{\arraystretch}{0.96}
\renewcommand{\multirowsetup}{\centering}
\begin{tabularx}{\linewidth}{|>{\centering\arraybackslash}p{2.20cm}|>{\raggedright\arraybackslash}p{2.45cm}|>{\raggedright\arraybackslash}X|>{\raggedright\arraybackslash}p{2.30cm}|}
\hline
\rowcolor{gray!20} \TableHeaderFirst{Tipo} & \TableHeader{Variable} & \TableHeader{Definición conceptual} & \TableHeader{Unidad / métrica} \\ \hline
\smash{\raisebox{-0.65\baselineskip}{\textbf{Independiente}}} & Arquitectura de captura & Configuración técnica usada para capturar y procesar la telemetría biótica. & Categórica: centralizada / perimetral. \\ \hline
\multirow[c]{4}{=}{\smash{\raisebox{-3.2\baselineskip}{\textbf{Dependiente}}}} & Cobertura efectiva de captura & Proporción del periodo observado con actividad válida y evidencia suficiente para reconstruir eventos de ordeño o telemetría de collar. & Porcentaje (\%). \\ \cline{2-4}
& Concordancia evento--campo & Coincidencia entre los eventos reconstruidos por el motor y la observación presencial, usando la identidad operativa Allflex. & Conteo, diferencia absoluta y razón. \\ \cline{2-4}
& Exposición de telemetría en claro & Magnitud del tráfico \ac{UDP} biótico transmitido sin autenticación ni cifrado observables sobre la \ac{LAN}. & Paquetes, duración y porcentaje. \\ \cline{2-4}
& Resiliencia operativa (\acs{MTTR}) & Tiempo promedio de recuperación del sistema de captura ante fallos eléctricos o de red. & Segundos (s); validación complementaria. \\ \hline
\end{tabularx}
\end{table}
\setlength{\tabcolsep}{6pt}

\subsection{Estrategia de muestreo y definición de la línea base}
La recolección combinó ejecuciones manuales controladas y programación automatizada. La fase exploratoria sirvió para reconocer firmas seriales, estabilidad del entorno y estructura del árbol de evidencia; luego se consolidó una rutina diaria con \textbf{nueve bloques}: dos de \textit{ordeño completo} (AM y PM) y siete de \textit{telemetría de collar}. Esta separación evita tratar como equivalentes sesiones que aportan evidencia distinta.


Cada bloque conserva tres fases: línea base inicial, captura principal y línea base final. El planificador \texttt{FincaScheduler} (v3.0) gestiona esa estructura: la Figura~\ref{fig:timeline_scheduler} muestra su distribución diaria y la Tabla~\ref{tab:calendarizacion_bloques} detalla cada bloque.

\begin{figure}[H]
\centering
\begin{tikzpicture}[x=0.645cm,y=1.05cm]
    \newcommand{\blockorde}[4]{%
        \fill[teal!18,draw=teal!60!black,rounded corners=2pt] (#1,0.18) rectangle (#2,1.18);
        \node[font=\fontsize{5}{5.4}\selectfont\bfseries] at ({(#1+#2)/2},0.86) {#3};
        \node[font=\fontsize{5}{5.4}\selectfont] at ({(#1+#2)/2},0.48) {#4};
    }
    \newcommand{\blocknormal}[4]{%
        \fill[cyan!16,draw=cyan!65!blue,rounded corners=2pt] (#1,0.18) rectangle (#2,1.18);
        \node[font=\fontsize{5}{5.4}\selectfont\bfseries] at ({(#1+#2)/2},0.86) {#3};
        \node[font=\fontsize{5}{5.4}\selectfont] at ({(#1+#2)/2},0.48) {#4};
    }

    \fill[gray!8] (0,-0.52) rectangle (24.35,1.36);
    \draw[thick,-Latex] (0,0) -- (24.45,0);
    \foreach \tick/\tag in {0/00:00,2/02:00,4/04:00,6/06:00,8/08:00,10/10:00,12/12:00,14/14:00,16/16:00,18/18:00,20/20:00,22/22:00,24/24:00}{
        \draw (\tick,0.09) -- (\tick,-0.09);
        \node[below,font=\scriptsize] at (\tick,-0.12) {\tag};
    }

    \blocknormal{0.00}{1.83}{normal\_7}{23:26}
    \blockorde{2.25}{4.25}{ordeño\_am}{02:15}
    \blocknormal{4.57}{6.95}{normal\_1}{04:34}
    \blocknormal{7.38}{9.78}{normal\_2}{07:23}
    \blocknormal{10.20}{12.25}{normal\_3}{10:12}
    \blockorde{13.00}{15.00}{ordeño\_pm}{13:00}
    \blocknormal{15.17}{17.50}{normal\_4}{15:10}
    \blocknormal{17.80}{20.13}{normal\_5}{17:48}
    \blocknormal{20.62}{23.00}{normal\_6}{20:37}
    \blocknormal{23.43}{24.25}{\scalebox{0.62}{normal\_7}}{23:26}

    \begin{scope}[yshift=-1.55cm]
        \fill[teal!18,draw=teal!60!black] (0,0.58) rectangle (0.42,0.86);
        \node[right,font=\scriptsize] at (0.50,0.72) {Ordeño completo (serial + antena + PCAP)};
        \fill[cyan!16,draw=cyan!65!blue] (0,0.12) rectangle (0.42,0.40);
        \node[right,font=\scriptsize] at (0.50,0.26) {Telemetría de collar (antena + PCAP secuencial)};
    \end{scope}
\end{tikzpicture}
\caption{Distribución temporal de los nueve bloques de captura automatizada a lo largo del ciclo diario (elaboración propia).}
\label{fig:timeline_scheduler}
\end{figure}

\vspace{0.25\baselineskip}

\begin{table}[H]
\centering
\caption{Programación diaria de bloques de captura automatizada (v3.0 del planificador).}
\label{tab:calendarizacion_bloques}
\scriptsize
\setlength{\tabcolsep}{2.5pt}
\renewcommand{\arraystretch}{1.08}
\begin{tabularx}{\linewidth}{|>{\centering\arraybackslash}p{2.45cm}|>{\centering\arraybackslash}p{1.55cm}|>{\centering\arraybackslash}p{2.05cm}|>{\raggedright\arraybackslash}X|}
\hline
\rowcolor{gray!20} \TableHeaderFirst{Bloque} & \TableHeader{Hora de inicio} & \TableHeader{Tipo} & \TableHeader{Finalidad operativa} \\ \hline
\texttt{ordeño\_am} & 02:15 & Ordeño completo & Captura simultánea de serial, antena y \ac{PCAP}; base para tiempos de ordeño, tandas y validación de eventos. \\ \hline
\texttt{normal\_1} & 04:34 & Telemetría de collar & Observación del canal biótico y tráfico de antena posterior al ordeño matutino. \\ \hline
\texttt{normal\_2} & 07:23 & Telemetría de collar & Seguimiento de actividad ordinaria del dominio de collar y exposición en red. \\ \hline
\texttt{normal\_3} & 10:12 & Telemetría de collar & Monitoreo matutino del canal biótico y del tráfico \ac{UDP} de la antena. \\ \hline
\texttt{ordeño\_pm} & 13:00 & Ordeño completo & Captura simultánea del proceso vespertino; referencia para contraste con observación presencial. \\ \hline
\texttt{normal\_4} & 15:10 & Telemetría de collar & Monitoreo posterior al ordeño vespertino y continuidad del canal biótico. \\ \hline
\texttt{normal\_5} & 17:48 & Telemetría de collar & Seguimiento de estabilidad vespertina de la \ac{LAN} y del dominio de collar. \\ \hline
\texttt{normal\_6} & 20:37 & Telemetría de collar & Seguimiento nocturno temprano de estabilidad de la \ac{LAN} y tráfico \ac{UDP}. \\ \hline
\texttt{normal\_7} & 23:26 & Telemetría de collar & Ventana de cierre diario y observación de fondo del dominio de collar. \\ \hline
\end{tabularx}
\end{table}


La línea base se acotó a la \textbf{fase inicial exploratoria}, comprendida entre el 23 de febrero y el 13 de marzo de 2026. Ese subconjunto permite fijar una referencia temprana de latencia, variación temporal de la latencia y exposición \ac{UDP}; por sí solo, no sostiene una medida final de $\eta$ ni una cadena de custodia serial--red completa.

\begin{table}[H]
\centering
\caption{Línea base de la fase inicial exploratoria (23/02/2026 al 13/03/2026).}
\label{tab:matriz_maestra_diagnostico}
\footnotesize
\renewcommand{\arraystretch}{1.3}
\setlength{\tabcolsep}{4pt}
\begin{tabularx}{\linewidth}{|>{\raggedright\arraybackslash}p{4.8cm}|>{\centering\arraybackslash}p{2.6cm}|>{\raggedright\arraybackslash}X|}
\hline
\rowcolor{gray!20} \TableHeaderFirst{Parámetro crítico} & \TableHeader{Valor} & \TableHeader{Alcance metodológico} \\ \hline
Cobertura temporal de la fase inicial & 8 visitas / 45 sesiones & Conjunto temprano usado para línea base de red y reconocimiento serial. \\ \hline
Latencia media línea base \acs{LAN} ($\bar{L}$) & 61,250 ms & Referencia temprana para los capítulos iniciales. \\ \hline
Variación temporal promedio de latencia & 2,756 ms & Referencia temprana de estabilidad del segmento. \\ \hline
Frecuencia de multidifusión inicial & 2,512 Hz & Indicador de exposición del dominio biótico en la fase inicial. \\ \hline
Fundamento para $\eta$ directa & No sostenible con esta fase & No hubo simultaneidad suficiente para sostener una $\eta$ inicial robusta. \\ \hline
Fundamento para cadena de custodia serial--red & No sostenible con esta fase & La evidencia concurrente madura aparece en fases posteriores. \\ \hline
\end{tabularx}
\setlength{\tabcolsep}{6pt}
\end{table}

\subsection{Procesamiento, validación y priorización de evidencia}
El motor \texttt{FincaDiag Modular} se usó como instrumento de reconstrucción y correlación, no como sustituto de la intervención perimetral. Separó sesiones de \textit{línea base}, \textit{telemetría de collar} y \textit{ordeño completo}; en estas últimas reconstruyó el flujo serial, agrupó tandas, contrastó eventos con observación presencial y verificó la coherencia entre serial, \acs{PCAP} y bitácora de campo. La Tabla~\ref{tab:clasificacion_visitas} resume la priorización aplicada para no mezclar evidencia exploratoria con datos aptos para contraste estadístico.

\begin{table}[H]
\centering
\caption{Síntesis funcional de las visitas usadas en el estudio.}
\label{tab:clasificacion_visitas}
\footnotesize
\renewcommand{\arraystretch}{1.22}
\begin{tabularx}{\linewidth}{|>{\raggedright\arraybackslash}p{3.0cm}|>{\raggedright\arraybackslash}p{3.1cm}|>{\raggedright\arraybackslash}X|>{\centering\arraybackslash}p{2.0cm}|}
\hline
\rowcolor{gray!20} \TableHeaderFirst{Fase} & \TableHeader{Período} & \TableHeader{Uso metodológico} & \TableHeader{Prioridad} \\ \hline
Exploratoria & 23/02--13/03 & Reconocimiento del canal, primeras firmas seriales y línea base temprana de red. & Media \\ \hline
Transición & 14/03--31/03 & Ajuste del analizador, mezcla de evidencia parcial y depuración del flujo de captura. & Media--Alta \\ \hline
Madura operativa & 01/04--05/04 & Primeras sesiones con cobertura más estable y evidencia serial--red suficiente. & Alta \\ \hline
Validada en campo & 06/04--09/04 & Comparación con observación presencial y calibración de eventos de ordeño. & Muy alta \\ \hline
Previa a la intervención consolidada & 10/04--10/05 & Muestra formal del Objetivo~4: 31 visitas y $n_{\text{pre}}=64$ sesiones con serial y PCAP simultáneos. & Muy alta \\ \hline
Posterior a la intervención & Desde 11/05 & Acumulación de $n_{\text{post}} \geq 30$ para el contraste final y validación operativa de la pasarela. & Muy alta \\ \hline
\end{tabularx}
\end{table}

La priorización siguió cuatro reglas: separar exploración de validación, no mezclar sesiones de collar con sesiones de ordeño completo, conservar los casos con $\eta=0\%$ para evitar sesgo de selección y fijar el corte previo a la intervención antes de analizar la condición instrumentada. Bajo este criterio, el tramo del 10 de abril al 10 de mayo de 2026 quedó definido como muestra previa a la intervención, mientras que las sesiones posteriores al 11 de mayo conforman la muestra posterior a la intervención.

Para cerrar la fase metodológica, la arquitectura se seleccionó mediante una matriz multicriterio de tipo Pugh. Se compararon aislamiento galvánico, gestión del dominio de red, independencia administrativa, no invasividad y determinismo del canal serial. Con esos criterios se adoptó una alternativa perimetral con captura aislada y observación controlada de red; la comparación técnica detallada se desarrolla en el Capítulo~\ref{chapter:chapter05}.

\section{Diseño metodológico de integración de la solución}
\label{sec:diseno-integracion-perimetral}
Seleccionada la alternativa de intervención, la pasarela perimetral quedó definida como una capa intermedia de captura y continuidad. No sustituye la lógica del fabricante; observa sin perturbar, reduce ruido operativo y conserva evidencia cuando el enlace de salida es inestable.

La integración se resume por planos de decisión en la Tabla~\ref{tab:criterios_integracion}. La comparación técnica detallada de componentes y configuraciones se desarrolla en el Capítulo~\ref{chapter:chapter05}.

\begin{table}[H]
\centering
\caption{Criterios metodológicos usados para integrar la pasarela perimetral.}
\label{tab:criterios_integracion}
\footnotesize
\renewcommand{\arraystretch}{1.22}
\begin{tabularx}{\linewidth}{|>{\raggedright\arraybackslash}p{2.7cm}|>{\raggedright\arraybackslash}X|>{\raggedright\arraybackslash}X|}
\hline
\rowcolor{gray!20} \TableHeaderFirst{Plano} & \TableHeader{Decisión de integración} & \TableHeader{Función metodológica} \\ \hline
Físico & Observación del enlace \ac{RS-232} mediante derivación no invasiva e interfaz con aislamiento galvánico. & Capturar evidencia sin cargar el bus ni modificar la operación del sistema comercial. \\ \hline
Red & Separación del dominio de telemetría \ac{UDP} con un conmutador administrable de Capa 2. & Reducir tráfico ajeno al flujo biótico y mantener capturas \ac{PCAP} comparables. \\ \hline
Lógico & Priorización de tramas relevantes, normalización de salida y cola local con reenvío. & Asegurar trazabilidad de eventos y continuidad ante conectividad rural variable. \\ \hline
Operativo & Paso de ejecuciones manuales a operación programada y servicio residente desde el 11 de mayo de 2026. & Definir la condición posterior a la intervención y medir resiliencia operativa mediante \acs{MTTR}. \\ \hline
\end{tabularx}
\end{table}

\section{Reevaluación y rediseño}
La reevaluación no cambió la arquitectura seleccionada; sirvió para revisar supuestos del diseño y ajustar componentes antes del contraste final. En esta etapa se separan dos asuntos: los ajustes ya ejecutados durante el despliegue y los criterios de validación que se completan con la muestra posterior a la intervención.


\subsection{Síntesis de ajustes de rediseño}
La reevaluación se concentró en tres ajustes detectados durante el despliegue. Para no convertir esta sección metodológica en una descripción de configuración, la Tabla~\ref{tab:ajustes_rediseno} resume el supuesto inicial, la evidencia que obligó a corregirlo y la decisión adoptada.

\begin{table}[H]
\centering
\caption{Ajustes incorporados durante la reevaluación de la solución.}
\label{tab:ajustes_rediseno}
\footnotesize
\renewcommand{\arraystretch}{1.22}
\begin{tabularx}{\linewidth}{|>{\raggedright\arraybackslash}p{3.0cm}|>{\raggedright\arraybackslash}X|>{\raggedright\arraybackslash}X|}
\hline
\rowcolor{gray!20} \TableHeaderFirst{Ajuste} & \TableHeader{Evidencia que lo motivó} & \TableHeader{Decisión metodológica} \\ \hline
Red perimetral & Tráfico administrativo, conflictos ARP y baja visibilidad de paquetes bióticos en capturas PCAP. & Sustituir el conmutador inicial por un GS305E con copia controlada de puerto y control de ráfagas. \\ \hline
Modelo operativo & El reinicio automático de un proceso por lote podía reprocesar sesiones ya publicadas. & Migrar a un servicio residente con vigilancia continua, detección de sesiones nuevas y drenado de la cola local. \\ \hline
Contrato de comunicación & La verificación estricta del canal seguro evidenció inconsistencias en los certificados de prueba. & Normalizar la autoridad certificadora y los certificados antes de la validación posterior. \\ \hline
\end{tabularx}
\end{table}

Estos ajustes no cambiaron el objetivo de la intervención; cerraron vacíos detectados al pasar del laboratorio al nodo de campo. La condición posterior a la intervención se define a partir de la activación de la configuración perimetral y del servicio residente en la Raspberry Pi, mientras que el contraste estadístico conserva la muestra previa a la intervención delimitada en esta metodología.

\subsection{Validación estadística y resiliencia operativa}
\label{subsec:protocolo_estadistico}

El contraste principal se fijó sobre la eficiencia de extracción ($\eta$) antes de disponer de los datos posteriores, para evitar selección retrospectiva de la prueba. La muestra pre-intervención corresponde a sesiones de ordeño completo entre el 10 de abril y el 10 de mayo de 2026, con serial y \ac{PCAP} simultáneos, parámetros estabilizados y conservación deliberada de los casos con $\eta = 0\%$; queda definida como $n_{\text{pre}} = 64$, $\bar{\eta}_{\text{pre}} = 4,49\%$ y $\eta_{\text{máx}} = 25\%$. La muestra posterior aplicó los mismos criterios para sesiones del 11 al 27 de mayo de 2026, alcanzando el cierre planificado con $n_{\text{post}} = 38$ y $\bar{\eta}_{\text{post}} = 22,80\%$, superando el umbral $n_{\text{post}} \geq 30$ requerido para el contraste Mann--Whitney~U.

La hipótesis nula establece $\eta_{\text{post}} \leq \eta_{\text{pre}}$ y la alternativa $\eta_{\text{post}} > \eta_{\text{pre}}$, con contraste unilateral y $\alpha = 0{,}05$. El protocolo sigue el Capítulo~\ref{chapter:chapter03}: Shapiro--Wilk para normalidad; si ambas muestras son normales, Levene y prueba t de Student o Welch según varianzas; si al menos una no es normal, la U de Mann--Whitney con estadístico $U$, valor $p$ y tamaño del efecto $r = Z/\sqrt{N}$. Dado el 54,7\% de valores nulos y la asimetría positiva de la muestra previa, la U de Mann--Whitney constituye el contraste esperado.

Como validación complementaria, la resiliencia operativa se evalúa con simulaciones controladas de interrupción eléctrica y caída del túnel de red, registrando el \ac{MTTR}. Se toma como referencia técnica un restablecimiento inferior a $60\,\text{s}$, sin que ese umbral sustituya el contraste estadístico principal.
